\section{Chapter 8 --- Conclusion, Limitations, and Future Work}\label{sec:chapter-8-conclusion-limitations-and-fut}

\section{Summary of the dissertation}\label{sec:summary-of-the-dissertation}

Chapter 7 explained the mechanisms behind the results. This chapter states what the dissertation has shown, where its evidence stops, and what should come next.

The dissertation set out to defend a single sentence: \textbf{supervision flows only from the verifiable region, yet the learned competence extends beyond it.} Every training example used anywhere in this work passed a deterministic compile--ground--execute--verify chain. During training-data harvest only, it also passed an agreement check against an independently computed oracle --- the oracle filtering candidates, never authoring them. Nothing the checks could not vouch for was learned from. Yet the systems trained this way did not stop at the projection line. The fully-local student holds its accuracy on precisely the hard-composite slice where its cloud teacher decays by 6--12 points, and it stays flat across the dev-to-test transition on which every cloud-briefed configuration falls steeply.

The system that carries this argument, CICADA, is built in three layers. A \emph{hard-verification layer} projects every property of the question-to-answer chain that can be projected onto structure or execution, and pins it with a deterministic check; the checks shrink the blind region layer by layer without ever removing it. An \emph{abstention layer} meets the unprojectable residue --- the purely semantic correspondence between plan and question --- with calibrated refusal, tested by three families of abstention traps built into the benchmark as first-class citizens. A \emph{learning layer} draws supervision only from what the first layer passes, and the competence it learns reaches back over part of what the first layer cannot see. The empirical summary is the five-system scoreboard on the held-out test set (n=2,285): teacher 69.76\% < Llama hybrid 75.97\% < Qwen hybrid 78.03\% < Llama fully-local 83.33\% < Qwen fully-local 85.65\%, with every pairing significant at McNemar $p<10^{-14}$. The fully-local champion exceeds its cloud teacher by +15.89 points (95\% CI [+14.07, +17.70]) and its own cloud-hybrid counterpart by +7.61 points (CI [+6.10, +9.13]), on a single local GPU at zero marginal API cost.

The four contributions announced in Chapter 1 were delivered as follows.

\textbf{Pillar 1, the bootstrap closed loop}, was instantiated three times. First, filtering the untrusted teacher's outputs into clean supervision (harvest yield 65.5\% oracle-correct on answerable questions). Second, a self-harvest round in which the student's own filtered generations exceeded the teacher's yield (76.6\% on Qwen, 78.1\% on Llama; these are harvest yields, a data-engine metric distinct from system accuracy). Third, the denoising distillation of the understanding stage that produced the fully-local stack.

\textbf{Pillar 2, gain decomposition}, separated the hard-verification layer's contribution from the learning layer's. The untrained zero-shot rung with full scaffolding reaches 70.4\% on the development set --- the scaffolding floor --- and the training ladder adds its gains on top of that floor. The v1$\rightarrow$v2.2 pipeline iteration showed that the projection boundary itself moves under engineering effort.

\textbf{Pillar 3, DPO near-miss toxicity}, is a mechanism-explained negative result with a cross-base dose curve. Preference suppression of verifier-rejected near-miss negatives degraded exactly the capability its pairs were richest in, mildly on one base and severely on the other, and a two-arm recipe-fix ablation failed in the direction the mechanism predicts.

\textbf{Pillar 4, benchmark methodology}, produced the 12,828-question benchmark with dual-implementation oracles at 99.88\% measured agreement and plan-level split isolation. Its development/confirmatory discipline survived an adversarial review --- including that review's findings against us, which were adopted and disclosed rather than argued away.

\section{Limitations}\label{sec:limitations}

A dissertation whose central instrument is a filter owes the reader an equally disciplined account of what its evidence does not cover. Each limitation below is stated with its factual bound: what the claim would need in order to widen, and how far the present evidence actually reaches. Where a concrete next experiment exists, it is named. None of these limitations was found by a reviewer first; each is carried forward from the point in Chapters 4--7 where it arose.

\textbf{A single self-built benchmark in a single domain.} All claims in this dissertation are claims about UK-procurement knowledge-graph question answering, on a benchmark we constructed ourselves. The architecture-level claims --- the three-layer decomposition and the value of verifier-gated distillation --- rest on within-domain ablations. No evidence here licenses their transfer to another schema or domain. The cross-schema probe of §8.3 is the named test.

\textbf{The hard-composite slice is a difficulty composite, not a holdout.} The slice on which the student's advantage is largest (Chapter 7) is defined as hard operators $\cup$ second-level rewrites $\cup$ abstention classes. Both training and test contain this class, with disjoint plans and novel surfaces. It measures robustness to difficulty, not out-of-distribution generalisation, and we have not named it "OOD" anywhere in this work --- the naming discipline is itself the limitation, honestly carried. The strict holdout the benchmark does contain is a single test-only template family of small n, reported here rather than in any headline.

\textbf{Compositional generalisation is unproven.} The pre-registered compositional probe (ood\_probe\_v1) was paused at its pilot by an explicit scope decision; the full \textasciitilde{}600-row evaluation was never run. What exists is disclosed in §7.7: a 10/10 novelty-gate self-audit, an honest weakening of one template's novelty claim, and a production-entry executability pilot showing asymmetric coverage (novel bridge compositions execute; the two chosen non-bridge compositions fail at compile time). No compositional claim is made on this evidence.

\textbf{The teacher--student comparison is asymmetric.} The teacher operates zero-shot in-domain while the students are fine-tuned on in-domain data. Student-beats-teacher results are therefore claims about the \emph{pipeline-plus-distillation recipe} --- that a bootstrapped local stack outperforms the deployed cloud configuration it was distilled from --- not claims that an 8B model outranks its teacher in model quality. The symmetric comparison is exactly what the compositional probe was designed to provide, since on its templates both teacher and student are zero-shot. The design is frozen, the pilot is complete, and the full run is deferred, not abandoned.

\textbf{The teacher is a single replicate on the final test set.} Provider-noise robustness lives on the development set, where three identical teacher runs establish a noise floor of 72.6\% $\pm$ 1.1 and the student beats each replicate separately. The final-test teacher pairing is labelled single-replicate wherever it appears. The +15.89-point gap exceeds the measured floor many times over, but the floor is the reason we can say so.

\textbf{Preference optimisation was unstable in this regime.} The DPO stage's gains are real but fragile: the same recipe that added +2.0 points on one base removed 5.4 on the other, and the recipe-level fixes failed. We report the champion configuration with its mechanism-level explanation, not as a recommended general recipe.

\textbf{The bootstrap converges early.} The second self-harvest round (r2) lifted harvest yields substantially (answerable 86.3\%, bridge 52\% $\rightarrow$ 75\%) yet delivered +0.4 points on the development set --- not significant, and adjudicated to a supplementary row under the pre-committed gate. The verifier-gated loop has a natural ceiling on this task; the first round captured essentially all of the transferable signal. Bootstrap claims in this dissertation are one-round claims.

\textbf{Development-derived cues touch 19 test rows.} Five unsupported-question cue strings were derived from development-set errors, and 19 of 2,285 test rows string-match them. The dual report bounds the effect: excluding those rows moves every headline by at most 0.21 points, with rankings and significance unchanged (the deltas are negative because the matched rows were mostly answered correctly by all systems). Two verbatim train--test duplicate questions, found by the adversarial review, are likewise disclosed and excluded from final statistics.

\textbf{Correctness is convention-relative, and surfaces are synthetic.} The oracles encode documented interpretive conventions; the 99.88\% dual-implementation agreement (14,752/14,770, residual 18) bounds but does not remove convention risk. The benchmark's question surfaces are LLM-generated across six style axes --- wider than template English, but not equal to real user language.

\section{Future work}\label{sec:future-work}

\textbf{Reinforcement learning from the verifier as reward.} The pipeline already contains what RLVR-style methods require: a dense, deterministic, non-learned reward source. Chapter 7's negative result argues for the direction --- additive, on-policy reinforcement (GRPO-style) rather than pairwise suppression of near-miss negatives --- and the r2 convergence result says what a new method must beat. One caveat is owed in advance. The verifier's blind region is exactly where reward hacking would live: a policy optimised against the verifier is optimised to \emph{pass checks}, and plans that are check-passing but semantically wrong receive positive reward by construction. Any RLVR follow-up must therefore carry the abstention layer and oracle-agreement audits into the reward loop, not just inherit them from the training-data pipeline.

\textbf{A cross-schema probe.} Porting the recipe --- three layers, projection-boundary engineering, verifier-gated self-harvest --- to an established text-to-SQL benchmark such as BIRD is the direct answer to the single-domain limitation. It tests whether the gain decomposition and the yield-surpassing self-harvest are schema-portable, on a benchmark this dissertation's authors did not build.

\textbf{The bootstrap convergence curve.} The three completed rounds --- teacher harvest, first self-harvest, second self-harvest --- sketch a curve whose yield keeps rising while its training gain stops. Two follow-ups suggest themselves. The first is to characterise the curve properly: does a third round's yield also translate to nothing, and does the point where gains stop move with model scale? The second is to test additive bridge re-consolidation, the lever the convergence analysis identified: raising the export cap on the bucket where the r2 harvest improved most, rather than adding rounds.

\textbf{Other procurement corpora.} The knowledge graph is built from OCDS-standard contract data, and the UK is one publisher among many. Rebuilding the graph over another country's OCDS releases would hold the schema family fixed while varying the data distribution, conventions, and language of the surfaces. This is a cheaper, intermediate generality test than a full schema change.

\textbf{The movable projection boundary as a research object.} The longer-term agenda turns the dissertation's mechanism into a question. For a given domain, which semantic properties can be projected into deterministic checks, at what engineering cost, and how does learned competence scale with the projected fraction? This work provides one measured data point --- a domain where partial projection plus abstention plus verifier-gated bootstrapping was enough to overtake the cloud teacher --- and the v1$\rightarrow$v2.2 iteration shows the boundary moves under engineering effort, in both directions. Treating partial verifiability as a general bootstrap principle, testable across domains where verification is inherently incomplete (data quality, code, scientific claims), turns the one-sentence thesis from a finding to be restated into a design law to be tested.

\section{Closing}\label{sec:closing}

The dissertation opened with a dilemma: the fluent guesser cannot be trusted with public data, and the mute executor cannot be asked a question. The answer defended here is neither to trust the language model nor to abandon it. Instead, we build the seam between the two out of checks wherever a check can reach, out of abstention wherever it cannot, and then let the checked region teach. The checks never verified everything, and they were never claimed to. They verified enough. The competence learned from that partial, filtered, deterministic signal reached beyond the region that produced it. Supervision flows only from the verifiable region, yet the learned competence extends beyond it.
